\section{Introduction}
\label{sec:intro}
\FigureArch

Politicians often answer the question they would have liked to be asked.
\citet{thomas-etal-2024-never} turned this into a classification task with a two-level taxonomy of response clarity and evasion, and SemEval-2026 Task~6 \citep{thomas-etal-2026-semeval} made it a shared task: Subtask~2 asks for one of nine evasion types, and Subtask~1 for one of three clarity levels that the nine types determine.

The competition was won by multi-call large language model (LLM) pipelines \citep{shi-etal-2026-teleai,peerachaidecho-rutherford-2026-chulanlp}, while fine-tuned encoders stopped at about 0.50 macro-F1 on Subtask~2 regardless of scale or ensembling \citep{thomas-etal-2026-semeval}.
A leaderboard does not say \emph{why} encoders stop there.
The candidate explanations make different predictions: the labels may be too noisy to learn (then annotators would disagree where models fail), training may optimise something other than the metric, which weighs all nine classes equally and accepts any annotator's label (then a better decision rule would help), the model may lack the knowledge needed to read an answer against its question (then a larger pretrained model would help under the same protocol), or the label meanings may not be learnable from 3,448 examples (then supplying definitions would help).

We test these explanations with a deliberately controlled single-pass classifier (Figure~\ref{fig:arch}): every change is made one at a time against a named control, on ten paired seeds, with predictions registered before each run and no choice made on the development set.
Our contributions are:
\begin{itemize}
  \item A DeBERTa-v3-large system (\S\ref{sec:encoder}) whose two real improvements (training length and the full journalist question) are separated from documented negative results: re-ranking, hierarchical models, rebalancing losses, per-class thresholds and model soups.
  \item Evidence that the remaining gap lies in the model rather than in label noise or the decision rule (\S\ref{sec:analysis}).
  \item A one-change test of the knowledge explanation (\S\ref{sec:llm}): replacing the backbone with Qwen3-8B-Base fine-tuned with LoRA \citep{qwen3,hu2022lora} improves both subtasks on all ten seeds and raises the ten-seed system from 0.405 to 0.543 (Subtask~2) and from 0.648 to 0.746 (Subtask~1).
  \item An analysis of where that gain falls, which refutes our own registered prediction (\S\ref{sec:analysis}).
\end{itemize}
Test labels were never released, so all comparisons are on the 308-item development set, with seed spreads, paired differences and bootstrap intervals throughout.
